% profiles/oregon-state-university.tex — Oregon State University
% ============================================================================
% source: https://graduate.oregonstate.edu/current-students/thesis-guide/formatting-thesis-or-dissertation
%         ("Formatting a Thesis or Dissertation", Office of Graduate Education —
%          the whole Thesis Guide, including the model pages it calls Figures
%          1-11; Figure 6 is the title page and Figure 7 the approval page)
% source: https://accessibility.oregonstate.edu/digital-accessibility
%         ("Electronic Documents Accessibility", which the Thesis Guide
%          incorporates by reference as a requirement)
% accessed: 2026-09-14
% checked:  texlib-thesis-profile-round
%
% VERIFIED and set here: geometry, spacing, title page, approval page, and the
% accessibility requirements including \thesisrequiretagging.
% NOT set: \thesissetchapteropening. OSU says "Each level 1 heading begins a new
% page" and states no deeper top margin for it.
%
% ONE REQUIREMENT THIS PROFILE CANNOT MEET, and it is not a small one: OSU puts
% the page number at the TOP RIGHT and leaves the pretext pages unnumbered.
%
%   "Pretext pages: Do not add page numbers to pretext pages."
%   "Page numbers must appear at the top right corner of pages, approximately 1
%    inch from the top edge of the page and at least 1 inch from the right edge
%    of the page. Page numbers must not invade any margins. There should be at
%    least one space between the page number and the first line of text."
%   "No headers or footers may be used. Footnotes are acceptable."
%
% This class puts the folio in the foot, and a profile owns page layouts, not
% the document's page style — the class builds its `normal' geometry with
% `includefoot' and there is no hook to move the number to a head. A document
% built against this profile is therefore OUT OF SPEC on pagination until
% either the class grows a hook or the author sets a page style themselves.
% Recorded here rather than half-done, because a wrong page number position is
% a reformat request, not a rejection.
%
% ALSO NOT ENCODABLE, recorded because each is a formatting rule OSU enforces:
% "No orphan lines may appear at the top or bottom of a page"; figures and
% tables must all be either inline or at chapter end, one system used
% consistently; "Choose high-contrast colors to differentiate lines, bars, or
% segments or use symbols with or without the color"; and the title "must be
% worded exactly the same throughout the document as it appears on the Abstract
% page, Title page and centered on page one".

\thesisinstitution{Oregon State University}

% "The left margin must be 1 inch unless printing and binding a personal or
% departmental copy then change to 1.5 inch. All other margins must be at least
% 1 inch, preferably 1.2 for top margin. Nothing may invade a margin. Every page
% must meet margin requirements."
%
% So: 1in on the left for the FILED copy — the 1.5in is for a bound personal
% copy and must not be in the electronic submission — and the top takes OSU's
% stated preference of 1.2in. Both 1in and 1.2in satisfy "at least 1 inch", so
% this profile follows the preference and an author who sets 1in is equally in
% spec; nothing here rests on the choice.
\thesissetgeometry{left=1in, right=1in, top=1.2in, bottom=1in}

% "Line spacing must be 1.5 or double, consistent throughout the document and
% matching which one you choose for the body of the thesis." A two-way choice
% with no stated preference, so this takes the class default. OSU names the
% single-spaced exceptions exactly: headings longer than one line, figure and
% table titles and legends, bibliographical and reference citations, direct
% quoted material, items listed within the body (optional), and "Where indicated
% in the pretext section".
%
% Type size is the author's within a range: "Use regular, unadorned print, 10-
% to 12-point size for text (headings may be 14-point only if all headings are
% 14-point)."
\thesissetspacing{double}

% --- Accessibility -----------------------------------------------------------
% OSU makes accessibility a filing requirement in the Thesis Guide itself, in
% the list of what the student is responsible for and again at the end:
%
%   "Ensuring that your document conforms to all requirements in this Thesis
%    Guide including those for electronic document accessibility."
%   "Your thesis or dissertation must meet OSU's Electronic Documents
%    Accessibility requirements."
%
% \thesisrequiretagging IS SET, and the reasoning should be visible because OSU
% never uses the word "tagged". It requires the filed document to be accessible,
% flatly; it recommends Word and says "Alternatively, PDF will be accepted if
% your thesis/dissertation contains complex formatting or technical
% requirements. Work closely with your major professor to ensure the digital
% accessibility of your PDF." A LaTeX thesis takes that PDF path by definition,
% and an untagged PDF cannot be an accessible one. If you read that as a
% requirement on Word documents only, drop the line.
%
% \thesisrequiredocumentlanguage and \thesisrequirepdfstandard are NOT set: OSU
% names neither the document language nor any PDF standard.
\thesisrequiretagging
\thesisaccessibilityrequirement{A filing requirement}
	{"Your thesis or dissertation must meet OSU's Electronic Documents
	 Accessibility requirements." The student is responsible for "Ensuring that
	 your document conforms to all requirements in this Thesis Guide including
	 those for electronic document accessibility."}
\thesisaccessibilityrequirement{Checked before submission}
	{"Microsoft Word is the recommended format for ensuring digital
	 accessibility. Run your document through Microsoft Word's accessibility
	 checker (Review>Check Accessibility), and resolve any issues identified.
	 Alternatively, PDF will be accepted if your thesis/dissertation contains
	 complex formatting or technical requirements. Work closely with your major
	 professor to ensure the digital accessibility of your PDF."}
\thesisaccessibilityrequirement{Alt text}
	{"Use the Alt Text functionality in Microsoft Word to add equivalent
	 alternative text to all images, figures, charts, and graphs."}
\thesisaccessibilityrequirement{Headings}
	{"Utilize Microsoft Word or LaTex Heading styles for effective and
	 accessible heading structures. Chapter names are Level 1 headings.
	 Subheadings of a chapter are Level 2 headings … Each level must look
	 different from the other levels. Headings of the same level must look the
	 same throughout the document." OSU is one of the few schools to name LaTeX
	 explicitly here.}
\thesisaccessibilityrequirement{Colour}
	{"Choose high-contrast colors to differentiate lines, bars, or segments or
	 use symbols with or without the color."}

% --- Profile-local metadata ---------------------------------------------------
% OSU distinguishes three things the class keeps as one or none:
%
%   * the DEFENCE date, which appears on the title page after "Presented" and
%     again in the approval page's opening sentence;
%   * the COMMENCEMENT date, which is the title page's last line — "Commencement
%     date is the June following the defense date, so if defense is after the
%     commencement ceremony it would be for the following year. Only month &
%     year, no date or it will be rejected." This profile takes it from
%     \graddate, the class's existing field;
%   * the DEPARTMENT, which is not the major: "Official major names and
%     department names can be found in the OSU Catalog. Some majors and
%     departments have the same name while others differ." \gradprogram carries
%     the MAJOR, and \oregonstatedepartment the department.
\newcommand{\thesis@osu@defense}{}
\newcommand{\oregonstatedefense}[1]{\renewcommand{\thesis@osu@defense}{#1}}
\newcommand{\thesis@osu@dept}{}
\newcommand{\oregonstatedepartment}[1]{\renewcommand{\thesis@osu@dept}{#1}}
\newcommand{\thesis@osu@needdefense}{%
	\ifx\thesis@osu@defense\@empty
		\ClassWarningNoLine{texlib-thesis}{OSU's title page and approval page both
			carry the^^J%
			DEFENCE date, which is not the commencement date. Set^^J%
			\string\oregonstatedefense{} (month spelled out, e.g. January 30,^^J%
			2027) or both pages will be filed incomplete}%
	\fi
}

% --- Title page --------------------------------------------------------------
% Figure 6 of the Thesis Guide, with its annotations, which are prescriptive:
%
%   <Title>                       "Title must match Abstract and page one title
%                                  exactly. Do not boldface the title."
%   by                            "Add two spaces after the title."
%   <Author>
%   A THESIS                      "Doctoral students may use 'A DISSERTATION'
%                                  instead of 'A THESIS' on Title Page,
%                                  Abstract, and Approval Pages."
%   submitted to
%   Oregon State University
%   in partial fulfillment of     "Follow division of this sentence (in partial
%   the requirements for the       fulfillment of...) exactly." — so these three
%   degree of                      line breaks are the specification, not a wrap
%   <Degree>                      "Spacing should be the same after your name,
%                                  'Oregon State University,' and your degree."
%   Presented <defence date>
%   Commencement <Month Year>
%
% "Titles longer than one line should be single-spaced."
%
% No page number: "Pretext pages: Do not add page numbers to pretext pages."
\renewcommand{\thesistitlepage}{%
	\loadgeometry{nopagenumber}%
	\thesis@osu@needdefense
	\begin{titlepage}%
		\thispagestyle{empty}\singlespacing\centering
		\vspace*{\fill}
		\thesis@tagtitle{Title}{\@title}\par
		\vspace{2\baselineskip}
		by\par
		\vspace{2\baselineskip}
		\@author\par
		\vspace{2\baselineskip}
		\MakeUppercase{A \thesis@DocNoun}\par
		\vspace{2\baselineskip}
		submitted to\par
		\vspace{2\baselineskip}
		Oregon State University\par
		\vspace{2\baselineskip}
		in partial fulfillment of\\
		the requirements for the\\
		degree of\par
		\vspace{2\baselineskip}
		\thesis@degree\par
		\vspace{2\baselineskip}
		Presented \thesis@osu@defense\par
		Commencement \thesis@date\par
		\vspace*{\fill}
	\end{titlepage}%
	\loadgeometry{normal}%
}

% --- Approval page ---------------------------------------------------------------
% Figure 7, "Standard Approval Page". Required, and filed WITHOUT signatures:
% "Your signature constitutes consent to have your document available for public
% reference in Valley Library, but the signatures on this page have been
% replaced with the ETD Submission Approval form", and the submission
% instruction is "Submit one copy of your thesis/dissertation, without
% signatures, electronically to ScholarsArchive." So the rules below are drawn
% empty on purpose, exactly as at Brown.
%
%   <Degree> thesis of <Author> presented on <defence date>.
%
%   APPROVED
%
%   ________________________
%   Major Professor representing <Major>
%
%   ________________________
%   Head of the Department of <Department>     ["If not part of a department,
%                                               please list the head/chair/dean
%                                               of the school or college."]
%   ________________________
%   Dean of Graduate Education
%
%   I understand that my thesis will become part of the permanent collection of
%   Oregon State University libraries. My signature below authorizes release of
%   my thesis to any reader upon request.
%
%   ________________________
%   <Author>, Author
%
% NOTE WHAT IS *NOT* ON THIS PAGE: no committee. OSU names four signatories —
% major professor, department head, Dean of Graduate Education, and the author —
% and the rest of the committee appears nowhere. \committeemember is therefore
% unused by this profile, deliberately; adding those lines would be inventing
% them.
%
% ALTERNATE WORDINGS the guide lists for cases this page does not generate, all
% replacing the first line or two: co-major professors ("Co-Major Professor,
% representing <Major>" twice), dual majors (one Co-Major Professor line per
% major), and MAIS ("Major Professor, representing Name of Major Area of
% Concentration" / "Co-Major Professor, representing Name of 2nd Major" /
% "Director of the Interdisciplinary Studies Program"). An author in one of
% those cases must override this page.
%
% Unnumbered, like every pretext page.
\renewcommand{\thesisapprovalpage}{%
	\clearpage\thispagestyle{empty}%
	\loadgeometry{nopagenumber}\singlespacing
	\thesis@osu@needdefense
	\ifx\thesis@osu@dept\@empty
		\ClassWarningNoLine{texlib-thesis}{OSU's approval page names the
			DEPARTMENT as well as^^J%
			the major, and they are not always the same. Set^^J%
			\string\oregonstatedepartment{}, or the head's line will read^^J%
			"Head of the Department of" with nothing after it}%
	\fi
	\thesis@tagtitle{H1}{{\color{white}\tiny Approval Page}}\par
	\vspace*{2\baselineskip}
	\noindent \thesis@degree{} \thesis@docnoun{} of \@author{} presented on
	\thesis@osu@defense.\par
	\vspace{3\baselineskip}
	{\centering APPROVED\par}%
	\vspace{2\baselineskip}
	\thesis@osu@sigline{Major Professor representing \thesis@program}%
	\thesis@osu@sigline{Head of the Department of \thesis@osu@dept}%
	\thesis@osu@sigline{Dean of Graduate Education}%
	\vspace{2\baselineskip}
	\noindent I understand that my \thesis@docnoun{} will become part of the
	permanent collection of Oregon State University libraries. My signature below
	authorizes release of my \thesis@docnoun{} to any reader upon request.\par
	\vspace{2\baselineskip}
	\thesis@osu@sigline{\@author, Author}%
	\vfill
	\clearpage\doublespacing
}

% One empty signature rule with its caption beneath, at the width Figure 7 draws
% it. Defined after the page that uses it; TeX resolves it at use time.
\newcommand{\thesis@osu@sigline}[1]{%
	\par\noindent
	\rule{4.2in}{0.4pt}\par
	% \noindent, or the caption picks up \parindent and sits a quad to the right
	% of the rule it belongs to. Figure 7 draws them flush.
	\noindent #1\par
	\vspace{2\baselineskip}
}
